% !TeX program = xelatex
% !TeX encoding = UTF-8
% !TeX spellcheck = en_US
% !BIB program = biber
%%
%% Practical work report: single-task GRPO on MolecularIQ.
%% Built on the JKU LaTeX report template by Michael Roland (MPL-2.0, see
%% LICENSE.template). Build with:  ./build.sh   (Linux/macOS)
%%                                 .\build.ps1  (Windows)
%%
%%%%%%%%%%%%%%%%%%%%%%%%%%%%%%%%%%%%%%%%%%%%%%%%%%%%%%%%%%%%%%%%%%%%%%%%%%%%%%%%

\documentclass[a4paper,oneside,10pt,ngerman,english]{scrartcl}

\usepackage[utf8]{inputenc}
\usepackage[techreport,fancyfonts,mathastext]{jkureport}

\usepackage{csquotes}
\usepackage[backend=biber,citestyle=numeric,sortcites=true,maxcitenames=2,style=ACM-Reference-Format]{biblatex}
\setcounter{biburlnumpenalty}{100}
\setcounter{biburllcpenalty}{100}
\setcounter{biburlucpenalty}{100}
\usepackage{todonotes}
\usepackage{import}
\usepackage{amsfonts}
\usepackage{booktabs}

\addbibresource{references.bib}

%%%%%%%%%%%%%%%%%%%%%%%%%%%%%%%%%%%%%%%%%%%%%%%%%%%%%%%%%%%%%%%%%%%%%%%%%%%%%%%%
%% Helper macros

%% \resultfig{path}: include a generated figure, degrading to a visible
%% placeholder if the pipeline has not been run, so the report always compiles.
\newcommand{\resultfig}[2][\linewidth]{%
    \IfFileExists{#2}{\includegraphics[width=#1]{#2}}{%
        \fbox{\parbox[c][3.2cm][c]{0.92\linewidth}{\centering\small\ttfamily
            figure not generated yet:\\#2\\[2pt]
            \normalfont run reproduce.sh report, then build.sh}}%
    }%
}

%% \tbd{...}: a value still to be filled in from a completed run. Deliberately
%% loud -- none of these may survive submission.
\newcommand{\tbd}[1]{\textcolor{red}{\textbf{[TBD: #1]}}}

%%%%%%%%%%%%%%%%%%%%%%%%%%%%%%%%%%%%%%%%%%%%%%%%%%%%%%%%%%%%%%%%%%%%%%%%%%%%%%%%

\begin{document}

\title{Single-Task GRPO on MolecularIQ}

\titleshort{Single-Task GRPO on MolecularIQ}

\subtitle{Technical Report\\%
    \usekomafont{subtitlesmall}%
    \hfill\\%
    Specialization, transfer, and what a verifiable reward actually buys\\
    \hfill\\%
}

%% FILL IN: your name and student e-mail address.
\author{%
    FIRSTNAME LASTNAME
    \affiliation{Institute for Machine Learning}
    \authornewline
    \authormail{firstname.lastname@students.jku.at}
    \authornewline
}

\revisionblock{%
    Practical work, supervised by SUPERVISOR NAME.
    All code, configurations and evaluation results are available in the
    accompanying repository; \texttt{scripts/reproduce.sh} regenerates every
    number and figure in this report.
}

\abstract{%
    We fine-tune Qwen2.5-0.5B-Instruct with Group Relative Policy Optimization
    (GRPO) separately on each of the three MolecularIQ task families --
    counting, indexing and constraint generation -- and evaluate all four
    resulting models on all three official benchmark splits. On the benchmark
    the picture is far weaker than in-distribution evaluation suggests: only
    two of three task-matched models beat the untrained baseline, counting
    degrades, and the strongest model on a given split is frequently not the
    one trained for it. A constant-answer control and answer-diversity
    statistics further show that our generated constraint questions are largely
    satisfiable without reasoning, and that GRPO amplified rather than resisted
    that shortcut -- the constraint-trained model places $91\%$ of its answers
    on a single molecule. We conclude that single-task GRPO at this scale buys
    little transferable capability, and that a verifiable reward needs a
    control measuring what the verifier accepts for free.
}

\keywords{reinforcement learning, GRPO, molecular reasoning, LLM evaluation, reward hacking}

\maketitle

\cleardoubleoddpage
\tableofcontents

%%%%%%%%%%%%%%%%%%%%%%%%%%%%%%%%%%%%%%%%%%%%%%%%%%%%%%%%%%%%%%%%%%%%%%%%%%%%%%%%

\cleardoubleoddpage

\section{Introduction}
\label{sec:introduction}

Large language models are increasingly asked to reason about molecules: to
count functional groups, locate atoms by index, or generate a structure
satisfying a set of numeric constraints. Unlike open-ended chemistry questions,
these tasks have a decisive property -- the answer can be checked mechanically.
MolecularIQ~\cite{bib:2024-moleculariq} exploits exactly this, pairing
generated questions over SMILES strings~\cite{bib:1988-weininger-smiles} with a
symbolic verifier built on RDKit~\cite{bib:rdkit}, so an answer is scored
objectively correct or incorrect rather than judged for plausibility.

A verifiable reward is what reinforcement learning needs. Group Relative Policy
Optimization (GRPO)~\cite{bib:2024-deepseekmath-grpo} samples a group of
completions per prompt, scores each with the verifier, and uses the
within-group advantage as its learning signal, dispensing with a separate value
network. This makes it practical to fine-tune a small model within a single GPU
allocation, and that is the setting explored here.

The natural next step from such a setup is multitask training. Before that is
worth attempting, a more basic question needs answering: \emph{how far does
training on one task alone get you, and does any of it transfer?} If a model
trained to count also improves at indexing, the tasks share structure that
multitask training could amplify. If each model improves only at its own task
-- or degrades elsewhere -- then multitask training is managing interference,
not harvesting synergy, and should be designed accordingly.

This report answers that question on the \emph{official} MolecularIQ benchmark.
That choice is deliberate and it is the central methodological commitment of
this work. It is tempting, and much cheaper, to evaluate on held-out questions
from one's own generator: they are unlimited, they match the training
distribution, and they produce encouraging numbers. We did exactly that
initially, and Section~\ref{sec:heldout-caveat} documents why we do not report
those numbers as results. Briefly, our generated evaluation set proved
systematically easier than the benchmark and, for constraint generation,
substantially reward-hackable -- to the point that a fixed string answering
every question scored as well as any trained model. Every headline result below
therefore comes from the official splits, and the generated held-out set
appears only as evidence about our own data.

We train three single-task models, one per task family, and evaluate all of
them on all three splits, forming a model~$\times$~task matrix whose diagonal
measures specialization and whose off-diagonal measures transfer. The
contributions are: (i) a single-task specialization-versus-transfer matrix for
MolecularIQ under GRPO, measured on the official benchmark; (ii) evidence that
in-distribution evaluation substantially inflates these results; and (iii) the
diagnosis of a reward-hackable flaw in generated constraint questions, together
with the control and diversity statistics that expose it.

\section{Background and Setup}
\label{sec:background}

\subsection{Tasks}

MolecularIQ groups questions into three families, which we treat as three
separate training tasks:

\begin{itemize}
\item \textbf{Counting} -- how many instances of a construct does a molecule
      contain (rings, aromatic rings, carbon atoms, hydrogen-bond donors,
      \dots)? The answer is an integer.
\item \textbf{Indexing} -- at which atom indices does a construct occur? The
      answer is a list of integers.
\item \textbf{Constraint generation} -- emit a SMILES string satisfying a set
      of property constraints (e.g.\ ``at least two rings and fewer than six
      rotatable bonds''). The answer is a molecule, checked by the verifier.
\end{itemize}

The first two are analytic: the molecule is given and must be read correctly.
The third is generative and admits many valid answers, which -- as
Section~\ref{sec:hacking} shows -- makes the difficulty of a question depend
entirely on how tightly its constraints were drawn.

\subsection{Model and training algorithm}

All runs fully fine-tune Qwen2.5-0.5B-Instruct~\cite{bib:2024-qwen25} (no
adapters) using the GRPO implementation in TRL~\cite{bib:trl} with vLLM for
rollout generation~\cite{bib:2023-vllm}. The reward is the official MolecularIQ
verifier score (binary, weight $1.0$) plus a small format-shaping term (weight
$0.2$) rewarding a parseable answer block; the shaping term exists only to stop
early training being starved of signal by malformed output. We use $16$
generations per prompt, a constant learning rate of $1\times10^{-6}$ after
warm-up, $500$ optimization steps, and $\beta=0$ (no KL penalty, hence no
reference model resident in memory). Full hyperparameters are in
\texttt{configs/*.yaml}.

\section{Method}
\label{sec:method}

\subsection{Training data}

Training questions are generated with \texttt{scripts/create\_datasets.py},
which draws molecules from the benchmark pool and emits questions in the
official schema: $20{,}000$ training questions per task, seed $42$, deduplicated.
For counting and indexing the fraction of questions whose answer is zero or
empty is capped at $25\%$, so a policy cannot reward-hack by always answering
``0''. This cap is \emph{not} applied to constraint generation, which
Section~\ref{sec:hacking} shows to be consequential.

\subsection{Evaluation}

All reported results are measured on the official
\texttt{ml-jku/moleculariq-v0.0} splits -- \texttt{single\_count},
\texttt{single\_index} and \texttt{single\_constraint\_generation} -- sampled at
$T=1.0$ with $n=3$ and scored as pass@3~\cite{bib:2021-chen-codex}, following
the benchmark paper's protocol. Scoring uses the official extraction function
and symbolic verifier, identical to the training reward.

We additionally maintain a held-out set of $500$ questions per task from our own
generator. It is used \emph{only} for checkpoint selection and as a diagnostic:
selecting checkpoints on the benchmark would bias the numbers we report. As
Section~\ref{sec:heldout-caveat} explains, it is not a valid measure of task
ability and does not appear in our results.

\subsection{Uncertainty and significance}

An accuracy computed on a finite question sample is an estimate, not a
constant, so every number carries an interval. For pass@$k$, which is strictly
$0/1$ per question, we use the Wilson score interval~\cite{bib:1927-wilson}; it
stays inside $[0,1]$ and remains well behaved near the extremes, which matters
here because several cells sit near zero. For mean scores we use a percentile
bootstrap over questions~\cite{bib:1993-efron-bootstrap}, which makes no
binomial assumption.

Comparisons against the baseline use a \emph{paired} bootstrap on the
difference, since both models answer the same questions in the same order;
pairing removes question difficulty as a nuisance factor and is markedly more
sensitive than checking whether two intervals overlap. For binary comparisons
an exact McNemar test~\cite{bib:1947-mcnemar} is also available. All estimators
live in \texttt{src/moleculariq\_grpo/stats.py} and are unit-tested against
closed-form values.

\subsection{Control baseline and diversity statistics}
\label{sec:control}

Because constraint generation admits many valid answers, a model can score
without reasoning by emitting one small molecule every time. To measure that
directly we add a \emph{constant-answer control}: a pseudo-model answering every
question with the fixed string \texttt{\{"smiles": "CC=O"\}}, requiring no model
and no GPU. It establishes the score reachable without reading the question, and
any result not clearing it is not evidence of task ability. We additionally
report answer-diversity statistics -- the number of distinct answers and the
share taken by the most common one -- for every evaluation.

\section{Results}
\label{sec:results}

\begin{figure}[!t]
\centering
\resultfig{figures/official-bars.pdf}
\caption{%
    Pass@3 on the official MolecularIQ v0.0 splits. Hatched gray is the
    constant-answer control, solid gray the untrained baseline, colours the
    single-task GRPO models. Error bars are $95\%$ Wilson intervals.
}\label{fig:official}
\end{figure}

\begin{table}[!b]
\centering
\caption{%
    Official benchmark results (pass@3). Rows are models, columns are splits;
    bold marks the task-matched cell of each column. The strongest model in a
    column is frequently \emph{not} the one trained for it.
    \tbd{confirm these values and add $95\%$ CIs and $p$-values from
    \texttt{results/matrix/plots/official/summary.csv}; the control row is
    outstanding.}
}\label{tab:official}
\small
\begin{tabularx}{\linewidth}{l>{\centering\arraybackslash}X>{\centering\arraybackslash}X>{\centering\arraybackslash}X}
\toprule
\textbf{Model} & \textbf{single\_count} & \textbf{single\_index} & \textbf{single\_constraint} \\
\midrule
Constant control    & \tbd{} & \tbd{} & \tbd{} \\
Qwen 0.5B untrained & 0.13 & 0.04 & 0.06 \\
GRPO count          & \textbf{0.07} & 0.04 & 0.12 \\
GRPO index          & 0.14 & \textbf{0.08} & 0.10 \\
GRPO constraint     & 0.07 & 0.00 & \textbf{0.14} \\
\bottomrule
\end{tabularx}
\end{table}

\subsection{Specialization is inconsistent}

The diagonal of Table~\ref{tab:official} does not support a clean
specialization story. Two of three task-matched models beat the untrained
baseline: indexing rises from $0.04$ to $0.08$ and constraint generation from
$0.06$ to $0.14$. The third moves the wrong way -- the count-trained model
scores $0.07$ on \texttt{single\_count} against a baseline of $0.13$, roughly
halving performance on the task it was trained for.

That inversion is the most informative single number in the study, because the
same model improves markedly on \emph{our own} counting questions
(Section~\ref{sec:heldout-caveat}). A policy that gets better on the
distribution it was trained on while getting worse on the benchmark has not
learned to count; it has learned our generator. Our questions are drawn with a
capped zero-answer fraction and a bounded SMILES length, making them narrower
than the official split -- shorter molecules, smaller answers, a different mix
of constructs. GRPO optimized for that narrower distribution at the benchmark's
expense.

The absolute magnitudes deserve emphasis too. Even the best cells sit between
$0.08$ and $0.14$ pass@3. A $0.5$B-parameter model remains far from competent at
these tasks after $500$ GRPO steps, and the differences we are discussing are a
few percentage points on splits of limited size. Whether individual comparisons
survive significance testing is exactly what the intervals in
Table~\ref{tab:official} must establish before any of them is leaned on.

\subsection{Transfer is absent}

\begin{figure}[!t]
\centering
\resultfig{figures/official-delta.pdf}
\caption{%
    Change in pass@3 relative to the untrained baseline on the official splits.
    Rows are models, columns are splits; outlined cells are the model's own
    task. Blue is improvement, red degradation.
}\label{fig:transfer}
\end{figure}

Off the diagonal, there is no evidence that training on one family helps
another in any systematic way, and some evidence it hurts. The
constraint-trained model drops to $0.00$ on indexing, below an untrained
baseline that was already only $0.04$. The count-trained model leaves indexing
unchanged at $0.04$.

Two off-diagonal cells are positive, and both are more awkward for a
specialization narrative than for a transfer one. The index-trained model is the
\emph{best} model on \texttt{single\_count} ($0.14$), beating both the baseline
and the count-trained model. Meanwhile the count-trained model improves
constraint generation ($0.12$ versus $0.06$). Neither is what one would predict
if each model had acquired its own task's procedure; both are consistent with
GRPO mainly teaching a general answer-formatting and answer-hygiene habit that
happens to help wherever the baseline was losing marks to malformed or rambling
output, and to hurt where it was not.

The interpretation we favour, then, is that at this scale and step budget GRPO
is not installing task-specific reasoning procedures at all. Counting requires
exhaustive enumeration over a molecular graph, indexing requires enumeration
plus positional bookkeeping, and constraint generation requires construction
rather than reading. Nothing in Table~\ref{tab:official} suggests any of these
procedures was learned; what moves is a few points of surface competence,
distributed almost arbitrarily across the three splits.

\subsection{GRPO amplifies a reward shortcut}
\label{sec:hacking}

The constraint results require a caveat that the benchmark numbers alone do not
reveal. Our generated constraint \emph{training} questions are, to a large
extent, satisfiable without reasoning, and GRPO found that shortcut.

The mechanism is in our generator. \texttt{\_build\_constraint} reads a property
value $v$ from a real anchor molecule and loosens it with a randomly chosen
operator. When $v=0$ -- common for constructs such as bridgehead atoms, E/Z
double bonds and R/S stereocenters, since most molecules have none -- every
operator branch yields a constraint that any molecule with zero of that property
satisfies: $=0$ holds, $\le 0+k$ holds, $<0+k$ holds, $\mathrm{range}[0,k]$
holds, and $>$ degenerates to $=0$ through a guard for $v<1$. A Monte-Carlo
replication of that function over $20{,}000$ draws confirms the consequence: at
$v=0$, $100\%$ of generated constraints are satisfied by a zero-valued molecule,
and $16.5\%$ collapse to a vacuous $\ge 0$ satisfied by every valid molecule. At
$v=1$ the figure is still $49.9\%$. The zero-answer cap that prevents this for
counting and indexing is not applied to constraint tasks
(\texttt{data.py:405}).

The consequence for training is visible in the model's outputs.
Table~\ref{tab:degenerate} reports answer diversity on the constraint task. The
untrained baseline already concentrates $43\%$ of its answers on a single small
molecule. After GRPO, the constraint-trained model concentrates \emph{$91\%$} of
its answers on one string, \texttt{\{"smiles": "CC(O)C"\}}. Training moved
diversity in the wrong direction: GRPO searched the space, found the degenerate
constant-output solution, and reinforced it, because that is what the reward
paid for. This is reward hacking, and it was not a training accident -- it was
the optimum of the objective we specified.

\begin{table}[!b]
\centering
\caption{%
    Answer diversity on the constraint task ($n=500$ generated questions).
    ``Top share'' is the fraction of questions answered with the single most
    common answer. Training increases concentration rather than reducing it.
}\label{tab:degenerate}
\small
\begin{tabularx}{\linewidth}{Xrrl}
\toprule
\textbf{Model} & \textbf{Distinct} & \textbf{Top share} & \textbf{Most common answer} \\
\midrule
Constant control      & 1   & $100\%$ & \texttt{\{"smiles": "CC=O"\}} \\
Qwen 0.5B untrained   & 62  & $43\%$  & \texttt{\{"smiles": "C(=O)C"\}} \\
GRPO count            & 70  & $43\%$  & \texttt{\{"smiles": "C(=O)C"\}} \\
GRPO index            & 86  & $35\%$  & \texttt{\{"smiles": "C(=O)C"\}} \\
GRPO constraint       & 44  & $\mathbf{91\%}$ & \texttt{\{"smiles": "CC(O)C"\}} \\
\bottomrule
\end{tabularx}
\end{table}

This does not by itself invalidate the constraint model's benchmark gain
($0.06 \rightarrow 0.14$): the official constraints are evidently far tighter
than ours, since the untrained baseline scores only $0.06$ there. But it means
the constraint-trained model should be regarded as a policy shaped by a
partly-degenerate objective, and its benchmark improvement as an incidental
by-product rather than as evidence of learned constraint satisfaction. Verifying
that would require retraining on repaired data, which we discuss in
Section~\ref{sec:limitations}.

\subsection{Why the generated held-out set is not reported}
\label{sec:heldout-caveat}

\begin{figure}[!b]
\centering
\resultfig[0.86\linewidth]{figures/heldout-bars.pdf}
\caption{%
    Our generated held-out questions, shown as a methodological caveat rather
    than as a result. Compare the constraint group against
    Table~\ref{tab:official}: the hatched constant-answer control -- a fixed
    string, no model -- scores $0.530$ here, statistically indistinguishable
    from every trained model.
}\label{fig:heldout}
\end{figure}

We initially evaluated on $500$ held-out questions per task from our own
generator, and those results are markedly more flattering than the benchmark.
Counting appears to rise from $0.120$ to $0.250$ and indexing from $0.000$ to
$0.250$ -- both large, both significant under a paired bootstrap, and both
substantially contradicted by the official splits, where counting in fact
degrades. We report them here only to document the size of the discrepancy.

Two independent problems make this set unrepresentative:

\begin{description}
\item[It is narrower than the benchmark.] Our questions are generated with a
      capped zero-answer fraction and a bounded SMILES length, and drawn from the
      same pool and generator as the training data. Improvements on it therefore
      partly measure fit to our own generator, which is precisely the failure the
      counting inversion exposes.
\item[Its constraint questions are largely trivial.] The constant-answer
      control scores $0.530$ ($95\%$ CI $[0.486, 0.574]$) on our constraint set,
      against $0.540$ for the untrained baseline and $0.552$ for the best trained
      model -- differences that are not significant. The same control scores
      exactly $0.000$ on counting and indexing, so this is not an artefact of the
      scoring pipeline but is specific to constraint generation, for the reason
      given in Section~\ref{sec:hacking}. In other words, the entire held-out
      constraint column is reachable with no model at all.
\end{description}

The general lesson generalizes beyond this project. An in-distribution held-out
set is a useful \emph{optimization} diagnostic -- it tells you whether training
fit the distribution it was given -- but it is not a measure of task ability, and
the gap between it and an independent benchmark is diagnostic in its own right.
Here that gap is what surfaced both the generator defect and the counting
inversion.

\section{Limitations}
\label{sec:limitations}

Several constraints on these conclusions should be stated plainly.

\emph{Effect sizes are small and split sizes are finite.} The best benchmark
cells sit at $0.08$--$0.14$ pass@3, and several comparisons differ by only a few
points. Until the confidence intervals in Table~\ref{tab:official} are filled
in, individual cell-to-cell comparisons should be treated as suggestive rather
than established. The qualitative findings we lean on -- the counting inversion,
the absence of systematic transfer, and the $91\%$ answer concentration -- are
large enough not to depend on that.

\emph{The constraint training data is defective.} As
Section~\ref{sec:hacking} documents, a large share of our constraint training
reward was obtainable without constraint reasoning. Repairing this means
extending the zero-answer cap to constraint tasks and rejecting vacuous
constraints, then regenerating the data and retraining -- work we did not have
the compute budget to redo. Until then the constraint-trained model is best
regarded as a reward-hacked artefact.

\emph{Single seed.} Each configuration was trained once. Our intervals capture
uncertainty from the finite evaluation sample, not from training stochasticity,
so a small difference between two trained models may reflect run-to-run
variance. Repeating training across seeds is the natural next step.

\emph{One model, one scale, one step budget.} All results are for a $0.5$B model
trained for $500$ steps. Both the absence of transfer and the reward hacking may
be scale- and budget-dependent; larger models or longer training could share more
structure across tasks, or resist degenerate solutions better. The negative
transfer result should not be extrapolated upward without evidence.

\emph{Single-construct training was not evaluated.} We also trained a variant
restricted to a single construct (counting aromatic rings). Because it is not
comparable to the three task-family models on the benchmark's task-level splits,
it is excluded here and left for future work.

\section{Conclusion}
\label{sec:conclusion}

Measured on the official MolecularIQ benchmark, single-task GRPO at this scale
buys very little. Two of three task-matched models improve on their own split
and one gets substantially worse; the strongest model on the counting split is
the one trained on indexing; and off-diagonal effects show no systematic
structure. Nothing in the results suggests that a task-specific reasoning
procedure -- graph enumeration, positional bookkeeping, molecular construction --
was actually acquired. What appears to move is a few points of general answer
hygiene, distributed almost arbitrarily across splits.

For the multitask work this practical leads into, that is still actionable, if
in a more sober direction than we expected. It suggests the three families do not
share enough procedural structure for gains to transfer implicitly, so a naive
mixture should be expected to face interference rather than synergy. It also
suggests that at $0.5$B and $500$ steps the more binding constraint may be
capability rather than task allocation: before tuning a task mixture, it is worth
establishing that any single task can be learned to a non-trivial level at all.
Designs worth trying next are those that manage interference explicitly --
balanced task sampling within a batch, per-task reward normalization so one
family's reward scale cannot dominate the group advantage -- evaluated on the
full matrix at every checkpoint, on the benchmark rather than in-distribution.

The methodological findings may be the more durable contribution. A benchmark
number is only as trustworthy as the control it is measured against. Our own
held-out set told a confident, encouraging, and wrong story: it reported counting
improving when the benchmark shows it degrading, and it scored a constant string
as highly as any trained model on constraint generation. Three inexpensive checks
caught this -- a second, independent evaluation source; an answer-diversity
statistic; and a constant-answer control. The last is worth stating as a general
rule: \emph{any reward defined by a verifier should be paired with a control that
measures what the verifier accepts for free.} Without it, the model that looks
best may simply be the one that has collapsed hardest, which is precisely what we
observed.

\printbibliography

\appendix

\section{Reproducing These Results}
\label{app:repro}

Every number and figure in this report is produced by the accompanying
repository. From a clean checkout with the pinned environment installed
(\texttt{scripts/setup\_leonardo.sh}):

\begin{lstlisting}[language=bash,caption={Full reproduction pipeline. Stages are idempotent and skip completed work.},label=lst:repro,float]
./scripts/reproduce.sh datasets  # generate training questions (seed 42)
./scripts/reproduce.sh train     # single-task GRPO runs
./scripts/reproduce.sh eval      # full model x split matrix
./scripts/reproduce.sh report    # figures + summary tables
pytest -q                        # unit tests for the statistics
\end{lstlisting}

Randomness is pinned throughout: dataset generation and training use seed $42$,
evaluation uses seed $0$. On the pinned software stack and the same GPU model
this reproduces the reported numbers exactly; across different GPU models,
floating-point non-determinism in the attention kernels can move individual
accuracies by a few tenths of a percent, well inside the reported confidence
intervals.

Per-cell numbers, confidence intervals, $p$-values and answer-diversity
statistics are written to \texttt{summary.csv} and \texttt{summary.md} in each
group's plot directory under \texttt{results/matrix/plots/}.

\cleardoubleoddpage

\end{document}
\endinput
